\documentclass[a4paper,12pt]{article}

% --- Pacchetti Necessari ---
\usepackage[utf8]{inputenc}
\usepackage[italian]{babel}
\usepackage[T1]{fontenc}
\usepackage{hyperref} % Per link cliccabili
\usepackage{xcolor}   % Per i colori del codice
\usepackage{listings} % Per visualizzare il codice SQL
\usepackage{booktabs} % Per tabelle eleganti
\usepackage{longtable} % Per tabelle su più pagine

% --- Configurazione Syntax Highlighting per SQL/PLSQL ---
\definecolor{codegreen}{rgb}{0,0.6,0}
\definecolor{codegray}{rgb}{0.5,0.5,0.5}
\definecolor{codepurple}{rgb}{0.58,0,0.82}
\definecolor{backcolour}{rgb}{0.95,0.95,0.92}

\lstset{
    language=SQL,
    backgroundcolor=\color{backcolour},   
    commentstyle=\color{codegreen},
    keywordstyle=\color{magenta},
    numberstyle=\tiny\color{codegray},
    stringstyle=\color{codepurple},
    basicstyle=\ttfamily\small,
    breaklines=true,
    captionpos=b,                    
    keepspaces=true,                 
    numbers=left,                    
    numbersep=5pt,                  
    showspaces=false,                
    showstringspaces=false,
    showtabs=false,                  
    tabsize=2
}

\title{Documentazione Tecnica Package PL/SQL \\ \textbf{BASEHTML}}
\author{Sviluppo Software}
\date{\today}

\begin{document}

\maketitle
\tableofcontents
\newpage

\section{Introduzione}
Il package \texttt{BASEHTML} è stato progettato per facilitare la generazione di interfacce web HTML direttamente da procedure PL/SQL in ambiente Oracle. Gestisce la struttura della pagina, i componenti di layout, i moduli di input e le tabelle.

\section{Variabili Globali}
Il package mantiene uno stato interno tramite le seguenti variabili:
\begin{itemize}
    \item \textbf{\texttt{v\_idSessione}}: (\texttt{NUMBER}) Memorizza l'identificativo della sessione utente. Valore di default: $-1$.
\end{itemize}

\section{Specifiche delle Procedure}

\subsection{Gestione Pagina}
\begin{itemize}
    \item \texttt{apriPagina(titolo, p\_idSessione)}: Apre il tag \texttt{html} e il \texttt{body}. Crea l'header con le procedure d'accesso e di navigazione del sito. Inoltre, gestisce il reindirizzamento in caso di sessione non valida
    \item \texttt{chiudiPagina}: Chiude i tag aperti della pagina e crea il footer.
\end{itemize}

\subsection{Elementi di testo}
I seguenti elementi servono per scrivere del testo all'interno del sito
\begin{longtable}{@{}p{5cm} p{8cm}@{}}
\toprule
\textbf{Procedura} & \textbf{Descrizione} \\ \midrule
\texttt{paragrafo}(testo, stile) & Inserisce un tag \texttt{p} con il testo e lo stile CSS fornito. \\
\texttt{h1}(testo, stile) & Inserisce un titolo \texttt{h1} con il testo e lo stile CSS fornito. \\
\texttt{collegamento}(testo, pagina, stile) & Inserisce un link \texttt{a} con il testo mostrato,la pagina a cui andare al click e lo stile CSS fornito. \\
\bottomrule
\end{longtable}

\newpage
\subsection{Layout e Contenitori}
Attraverso queste procedure è possibile definire la struttura della pagina.

\begin{longtable}{@{}p{5cm} p{8cm}@{}}
\toprule
\textbf{Procedura} & \textbf{Descrizione} \\ \midrule
\texttt{apriDiv}(id, stile) & Apre un tag \texttt{div}. Accetta \texttt{id} con flex, lista e griglia che forniscono stili predefiniti (sia alla div stessa che agli elementi interni ad essa) e \texttt{stile} per ritocchi personalizzati di CSS. \\
\texttt{chiudiDiv} & Chiude il tag \texttt{div}. \\
\texttt{apriPopup}(id) & Apre un popup \texttt{dialog} con lo stile dell'ID fornito. \\
\texttt{chiudiPopup} & Chiude il tag \texttt{dialog} \\
\bottomrule
\end{longtable}

\subsection{Form e Input}
Queste procedure permettono la creazione di moduli interattivi.

\begin{longtable}{@{}p{5cm} p{8cm}@{}}
\toprule
\textbf{Procedura} & \textbf{Descrizione} \\ \midrule
\texttt{apriModulo}(id, action) & Apre un tag \texttt{form}. Accetta \texttt{id} con modulo per l'incorniciamento del form e \texttt{action} indica la pagina a cui andare e inviare i dati. \\
\texttt{chiudiModulo} & Chiude il tag \texttt{form}. \\
\bottomrule
\end{longtable}

All'interno di un form, ogni campo deve avere un nome per poter essere passato alla pagina di destinazione all'invio del form
Per inserire campi in un form, ci sono diverse opzioni:
\subsubsection{Input}
Inserisce un input standard nel form

\begin{lstlisting}[caption={firma della procedura InserisciInput}]
PROCEDURE inserisciInput(
    id           IN VARCHAR2,
    tipo         IN VARCHAR2 DEFAULT 'text',
    nome         IN VARCHAR2,
    valore       IN VARCHAR2 DEFAULT NULL,
    placeholder  IN VARCHAR2 DEFAULT NULL,
    obbligatorio IN BOOLEAN  DEFAULT false
);
\end{lstlisting}

\paragraph{Parametri di \texttt{inserisciInput}:}
\begin{itemize}
    \item \textbf{id}: indica il testo del label relativo all'input, se lasciato vuoto, non si crea nessun label
    \item \textbf{tipo}: Indica il tipo di input.
        \subitem text = tipo di default, una casella di testo di una riga
        \subitem password = una casella di testo che nasconde i caratteri al suo interno
        \subitem email = una casella di testo, richiede una @ e un . dopo la @
        \subitem checkbox = indica un elemento di una lista con una sola opzione valida
        \subitem radio = indica un elemento di una lista con più opzioni valide
        \subitem date = input di tipo data, mostra il calendario e richiede una data
        \subitem hidden = l'input non viene mostrato all'utente, utile per passare parametri nascossti come l'id della sessione
    \item \textbf{nome}: Il nome che verrà usato per recuperare il valore del campo, è quello che appare nell'URL della pagina all'invio del form.
    \item \textbf{valore}: valore di default del campo, prima che l'utente lo modifichi (per input di tipo hidden è il valore che vogliamo passare)
    \item \textbf{placeholder}: valore mostrato in caso di campo vuoto. solitamente è un'istruzione come: inserisci email, inserisci password. sparisce appena si compila il campo
    \item \textbf{obbligatorio}: Se impostato a \texttt{true}, aggiunge l'attributo \texttt{required}, e non si può inviare il form fino a che il campo non è stato compilato.
\end{itemize}

\subsubsection{Menu A Tendina}
inserisce un menù a tendina nel form
\begin{itemize}
    \item \textbf{apriMenuTendina(id, nome, stile)}: Apre il tag \texttt{select}, ID indica il testo del label relativo all'input, se lasciato vuoto, non si crea nessun label, Nome è Il nome che verrà usato per recuperare il valore del campo, è quello che appare nell'URL della pagina all'invio del form e Stile è lo stile CSS aggiungibile manualmente.
    \item \textbf{chiudiMenuTendina}: Chiude il tag \texttt{select}
    \item \textbf{tendinaOption(opzione, valore, selezionato)}: aggiunge un opzione al menù a tendina. opzione è il nome mostrato all'utente, valore è il nome mostrato alla macchina(=opzione se valore è omesso) e selezionato indica l'opzione da selezionare al caricamento della pagina

\end{itemize}

\subsubsection{Bottone}
Ogni form richiede un bottone per inviare il form, ma può essere usato anche altrove per richiamare funzioni

\begin{longtable}{@{}p{5cm} p{8cm}@{}}
\toprule
\textbf{Procedura} & \textbf{Descrizione} \\ \midrule
\texttt{Bottone}(testo, onClick) & Inserisce un tag \texttt{button} con il testo indicato e la funzione da chiamare. per i form, omettere il paramtro onClick.  \\
\bottomrule
\end{longtable}

\subsubsection{TextArea}
Utilizzabile quando si vuole un testo in un riquadro di dimensione aggiustabile, comunemente utilizzato nei form per inserire input più lunghi di una riga.

\begin{longtable}{@{}p{5cm} p{8cm}@{}}
\toprule
\textbf{Procedura} & \textbf{Descrizione} \\ \midrule
\texttt{inserisciTextArea}(testo, nome, modificabile) & Inserisce un tag \texttt{TextArea} con il testo indicato. Nome permette di richiamare la textarea dopo un form, mentre il parametro modificabile indica se l'utente può modificare o meno il campo(default true) \\
\bottomrule
\end{longtable}

\subsection{Tabella}
I seguenti tag servono per creare una tabella

\begin{longtable}{@{}p{5cm} p{8cm}@{}}
\toprule
\textbf{Procedura} & \textbf{Descrizione} \\ \midrule
\texttt{apriTabella}(id, stile) & Inserisce un tag \texttt{table} con lo stile CSS indicato (id per tabelle non definiti)\\
\texttt{chiudiTabella} & chiude il tag \texttt{table}\\
\texttt{inizioRiga} & Inserisce un tag \texttt{tr} che inizia la riga della tabella\\
\texttt{fineRiga} & Inserisce un tag \texttt{tr} che finisce la riga della tabella\\
\texttt{inserisciIntestazione}(testo) & Inserisce un tag \texttt{th} che inserisce nella riga una cella di tipo intestazione con il testo indicato\\
\texttt{inserisciCella}(testo) & Inserisce un tag \texttt{td} che inserisce nella riga una cella con il testo indicato\\
\texttt{apriCella} & Apre un tag \texttt{td} per poter mettere nella cella elementi diversi dal testo\\
\texttt{chiudiCella} & Chiude un tag \texttt{td} per chiudere una cella aperta con apriCella\\


\bottomrule
\end{longtable}

\section{Utility di Sistema}
\begin{itemize}
    \item \texttt{aggiungi\_stile(stile)}: Permette di inserire blocchi CSS custom.
    \item \texttt{aggiungi\_script(script)}: Inserisce codice JavaScript nella pagina.
    \item \texttt{redirect(url)}: Forza un reindirizzamento della pagina tramite browser.
\end{itemize}

\section{Esempio di Utilizzo}
Di seguito un esempio di codice PL/SQL per generare una pagina più semplice possibile utilizzando il package \texttt{BASEHTML}:

\begin{lstlisting}
create or replace procedure esempio( p_idSessione IN NUMBER DEFAULT -1) is
BEGIN
   BASEHTML.apriPagina('Login Sistema', idSessione);
   
   BASEHTML.chiudiPagina;
END;
end home;
/
GRANT EXECUTE ON esempio TO anonymous;
\end{lstlisting}
mettendo il grant della pagina in fondo alla pagina stessa, si può concedere l'autorizzazione alla pagina quando la si crea

\end{document}